\chapter{Đặc tả yêu cầu hệ thống}\label{chap:requirements}
\section{Danh sách use case}
Các yêu cầu được chia thành yêu cầu chức năng (FR), yêu cầu phi chức năng (NFR) và quy tắc nghiệp vụ (BR). Mỗi use case (UC) mô tả một mục tiêu của người sử dụng; các thao tác con được đánh số để thuận tiện đối chiếu với kế hoạch kiểm thử.

\paragraph{DEC07 -- Xác thực và quản lý phiên.}
iCareX sử dụng cookie HttpOnly để duy trì đăng nhập và lưu phiên trong cơ sở dữ liệu. Mỗi yêu cầu được kiểm tra theo tài khoản, thời hạn và trạng thái thu hồi của phiên. Các thao tác thay đổi dữ liệu phải có CSRF token hợp lệ. Khi tải lại trang, ứng dụng gọi \code{/api/v1/auth/me} để khôi phục thông tin tài khoản nếu phiên còn hiệu lực. Đăng xuất thu hồi phiên tương ứng và xóa cookie trên trình duyệt~\cite{csrf}.

\begin{smalltable}\begin{longtable}{L{1.3cm}L{3.5cm}L{5.7cm}L{1.8cm}L{2cm}}
\caption{Danh mục use case và phạm vi}\label{tab:uc-list}\\
\toprule ID & Tên & Phân rã thao tác & FR & Ưu tiên\\\midrule\endfirsthead
\toprule ID & Tên & Phân rã thao tác & FR & Ưu tiên\\\midrule\endhead
UC01 & Tài khoản và phiên & .1 Đăng ký; .2 Đăng nhập; .3 Đăng xuất & FR01 & Cao\\
UC02 & Quản lý xe & .1 Thêm; .2 Sửa; .3 Archive/restore & FR02 & Cao\\
UC03 & Quản lý ODO & .1 Thêm manual; .2 Sửa/xóa manual & FR03 & Cao\\
UC04 & Nhiên liệu & .1 Thêm; .2 Sửa/xóa & FR04 & Cao\\
UC05 & Dịch vụ / completion & .1 Thêm nhiều dòng; .2 Sửa/xóa và tính lại & FR05,08 & Cao\\
UC06 & Kế hoạch & .1 Tạo/sửa; .2 Pause/resume/xóa có điều kiện & FR06 & Cao\\
UC07 & Trạng thái và nhắc & .1 Xem hạn; .2 Đọc inbox & FR07,09 & Cao\\
UC08 & Lịch sử & .1 Tra cứu, lọc, mở nguồn & FR10 & Cao\\
UC09 & Dashboard / báo cáo & .1 Tổng quan; .2 Báo cáo & FR11 & Cao\\
UC10 & Hồ sơ / cài đặt & .1 Hồ sơ; .2 Mật khẩu; .3 Cài đặt & FR01,12 & Cao\\
UC11 & Catalog & .1 Đọc cascade; .2 Quản trị metadata & FR02 mở rộng & Bổ sung\\\bottomrule
\end{longtable}\end{smalltable}

\section{Sơ đồ use case tổng quát và phân hệ}
\diagram{uc-overall.tex}{Use case tổng thể: actor ngoài ranh giới, các chức năng trong iCareX}{uc-overall}
Khách thực hiện đăng ký và đăng nhập. Chủ xe sử dụng các chức năng quản lý dữ liệu cá nhân, còn quản trị danh mục duy trì thông tin loại xe, hãng và dòng xe.
\diagram{uc-vehicle.tex}{Use case xe/ODO; kiểm tra chuỗi là hành vi bắt buộc khi ghi reading}{uc-vehicle}
\diagram{uc-records.tex}{Use case hoạt động: liên kết completion mở rộng việc ghi dịch vụ}{uc-records}
Liên kết completion là tùy chọn tại extension point chọn hạng mục Maintenance; mũi tên extend đi từ hành vi tùy chọn về ghi dịch vụ. Khi không chọn plan, dịch vụ vẫn ghi lịch sử và chi phí. Transaction là hành vi nội bộ bắt buộc của việc ghi dịch vụ trong thiết kế đích.
\diagram{uc-plans.tex}{Use case kế hoạch/nhắc: đọc nhắc độc lập với ghi việc đã làm}{uc-plans}

\section{Yêu cầu chức năng}
\begin{smalltable}\begin{longtable}{L{1.2cm}L{3.1cm}L{10.3cm}}
\caption{Yêu cầu chức năng và điều kiện hoàn thành}\label{tab:fr}\\
\toprule Mã & Nhóm & Hành vi yêu cầu\\\midrule\endfirsthead
\toprule Mã & Nhóm & Hành vi yêu cầu\\\midrule\endhead
FR01 & Tài khoản & Đăng ký không tự login; phiên riêng; logout thu hồi; xem/sửa tên, email chỉ đọc; đổi password cần password cũ và thu hồi phiên.\\
FR02 & Phương tiện & Thêm/xem/sửa; archive/restore giữ lịch sử. Chọn model thuộc brand/type active; plate unique theo owner kể cả archive.\\
FR03 & ODO & Thêm/sửa/xóa manual; giữ chuỗi theo ngày; reading từ nguồn phải sửa qua nguồn; current ODO không lùi khi nhập quá khứ.\\
FR04 & Nhiên liệu & Ngày, ODO, lít, tiền, ghi chú; tạo có RequestId; sửa/xóa có version; nguồn và reading cùng transaction.\\
FR05 & Dịch vụ & Ít nhất một dòng Repair/Maintenance, tổng dương; mỗi Maintenance có thể liên kết plan cùng xe. Lỗi một dòng rollback cả dịch vụ.\\
FR06 & Kế hoạch & Time/Distance/Either; baseline có ý nghĩa rõ, chu kỳ/ngưỡng; pause/resume; chỉ xóa nếu chưa completion.\\
FR07 & Trạng thái & Unknown/NotDue/Soon/Due/Overdue; reason, ngày/km còn lại, ngày cập nhật ODO và cờ dữ liệu cũ.\\
FR08 & Completion & Chỉ dòng Maintenance liên kết mới chuyển kỳ; sửa/xóa phải tìm latest hợp lệ còn lại; không cộng chi phí hai lần.\\
FR09 & Notification & Inbox DB, đọc, badge, Superseded; unique plan/kỳ/mức; worker startup/định kỳ; browser permission fallback.\\
FR10 & Lịch sử & Lọc xe/ngày/loại/từ khóa, phân trang ổn định; một nguồn không thành hai hoạt động do reading tự sinh.\\
FR11 & Tổng hợp & Dashboard và báo cáo chung filter; fuel amount + service item amount; nhóm tháng/xe/loại; số liệu đối chiếu được.\\
FR12 & Cài đặt / UX & Nhắc, ngưỡng ODO cũ, trạng thái loading/error/empty/conflict; responsive và thao tác bàn phím.\\\bottomrule
\end{longtable}\end{smalltable}

\subsection{Tiền điều kiện và hậu điều kiện chung}
Các thao tác dữ liệu lấy account từ phiên server và kiểm tra quan hệ owner/vehicle; không dùng accountId trong payload làm nguồn phân quyền. Xe archived chỉ đọc. Query không thay dữ liệu nghiệp vụ. Thành công mutation chỉ được báo sau commit; lỗi phải giữ bản ghi nhất quán. Validation phía UI hỗ trợ trải nghiệm, không thay kiểm tra server. Version cũ trả 409 để tải lại; form đang nhập được giữ khi lỗi mạng, không báo lưu thành công từ dữ liệu trong memory.

\subsection{Đặc tả chi tiết từng thao tác}
Mỗi bảng use case trình bày tác nhân, điều kiện bắt đầu, luồng xử lý, ngoại lệ và kết quả mong đợi. Các mã AT liên kết use case với tiêu chí nghiệm thu tương ứng.
\subsubsection{UC01.1 -- Đăng ký tài khoản}\label{uc:UC01.1}
\begin{smalltable}\begin{longtable}{L{3.2cm}L{11.8cm}}
\caption{UC01.1 -- Đăng ký tài khoản}\label{tab:UC01.1}\\
\toprule Trường & Đặc tả\\\midrule\endfirsthead
\toprule Trường & Đặc tả\\\midrule\endhead
\textbf{ID / tác nhân / ưu tiên} & UC01.1; Khách; Ưu tiên cao.\\
\textbf{Yêu cầu / kích hoạt} & FR01; Chọn Đăng ký.\\
\textbf{Tiền điều kiện} & Có kết nối mạng; CSRF token hợp lệ.\\
\textbf{Luồng chính} & \begin{enumerate}[leftmargin=*,nosep]
\item Nhập tên, email, password, xác nhận password và phone tùy chọn; UI kiểm tra trường.
\item Server chuẩn hóa email/tên, kiểm tra email và policy password.
\item Trong transaction, tạo Identity user, role User, Account và UserSettings.
\item Commit rồi trả 201 AccountResponse; không trả hash/password.
\item Hiển thị đăng ký thành công và chuyển người dùng sang đăng nhập.
\end{enumerate}\\
\textbf{Luồng thay thế} & Hủy form không ghi dữ liệu. Xác nhận password là kiểm tra UX; server vẫn tự enforce policy.\\
\textbf{Ngoại lệ} & Email trùng: 409; trường sai/CSRF sai: 400. DB lỗi không để account thiếu settings.\\
\textbf{Hậu điều kiện / chấp nhận} & Có một bộ user/profile/settings nhất quán; chưa có session login.\\
\textbf{Kiểm thử liên quan} & AT01, AT17.\\
\bottomrule\end{longtable}\end{smalltable}
\subsubsection{UC01.2 -- Đăng nhập và khôi phục phiên hợp lệ}\label{uc:UC01.2}
\begin{smalltable}\begin{longtable}{L{3.2cm}L{11.8cm}}
\caption{UC01.2 -- Đăng nhập và khôi phục phiên hợp lệ}\label{tab:UC01.2}\\
\toprule Trường & Đặc tả\\\midrule\endfirsthead
\toprule Trường & Đặc tả\\\midrule\endhead
\textbf{ID / tác nhân / ưu tiên} & UC01.2; Khách / chủ xe; Ưu tiên cao.\\
\textbf{Yêu cầu / kích hoạt} & FR01; Gửi biểu mẫu đăng nhập.\\
\textbf{Tiền điều kiện} & Tài khoản đã đăng ký; auth service và DB truy cập được.\\
\textbf{Luồng chính} & \begin{enumerate}[leftmargin=*,nosep]
\item Browser lấy CSRF token rồi gửi credentials qua kênh HTTPS khi triển khai.
\item Server kiểm tra CSRF, tìm user chuẩn hóa, kiểm tra lockout và password.
\item Tạo AuthSession trong DB với thời hạn 8 giờ; reset failed count khi thành công.
\item Controller phát cookie HttpOnly chứa identity/role/sid; trả thông tin tài khoản.
\item Khi khởi tạo UI, gọi /auth/me để lấy account nếu cookie/session còn hợp lệ.
\end{enumerate}\\
\textbf{Luồng thay thế} & Khi tải lại trang, trình duyệt gửi cookie phiên. Ứng dụng lấy thông tin qua /auth/me nếu phiên trong DB còn hiệu lực; nếu phiên hết hạn hoặc bị thu hồi, người dùng được yêu cầu đăng nhập lại.\\
\textbf{Ngoại lệ} & Thông tin sai/locked: 401 chung; CSRF: 400; DB: 503. Không phân biệt email tồn tại bằng thông báo login.\\
\textbf{Hậu điều kiện / chấp nhận} & Chỉ phiên chưa revoke/chưa hết hạn truy cập dữ liệu riêng; UI nhận account.\\
\textbf{Kiểm thử liên quan} & AT01.\\
\bottomrule\end{longtable}\end{smalltable}
\subsubsection{UC01.3 -- Đăng xuất}\label{uc:UC01.3}
\begin{smalltable}\begin{longtable}{L{3.2cm}L{11.8cm}}
\caption{UC01.3 -- Đăng xuất}\label{tab:UC01.3}\\
\toprule Trường & Đặc tả\\\midrule\endfirsthead
\toprule Trường & Đặc tả\\\midrule\endhead
\textbf{ID / tác nhân / ưu tiên} & UC01.3; Chủ xe; Ưu tiên cao.\\
\textbf{Yêu cầu / kích hoạt} & FR01; Chọn Đăng xuất.\\
\textbf{Tiền điều kiện} & Phiên hợp lệ; CSRF token.\\
\textbf{Luồng chính} & \begin{enumerate}[leftmargin=*,nosep]
\item UI gửi POST logout với cookie và CSRF header.
\item Server đọc account/sid từ claim và revoke đúng phiên trong DB.
\item Controller xóa cookie và trả 204.
\item UI xóa account, query cache và dữ liệu nhạy cảm đang mở.
\item Request tiếp theo với phiên cũ phải bị từ chối.
\end{enumerate}\\
\textbf{Luồng thay thế} & Nếu phiên đã hết hạn, UI chuyển trạng thái chưa đăng nhập sau 401. Revoke đã thực hiện có thể xử lý idempotent.\\
\textbf{Ngoại lệ} & Mạng lỗi: không thông báo đã revoke; cho thử lại. CSRF sai: không đổi session.\\
\textbf{Hậu điều kiện / chấp nhận} & Phiên tương ứng không dùng được; cache không còn dữ liệu account cũ.\\
\textbf{Kiểm thử liên quan} & AT01, AT02.\\
\bottomrule\end{longtable}\end{smalltable}
\subsubsection{UC02.1 -- Thêm phương tiện}\label{uc:UC02.1}
\begin{smalltable}\begin{longtable}{L{3.2cm}L{11.8cm}}
\caption{UC02.1 -- Thêm phương tiện}\label{tab:UC02.1}\\
\toprule Trường & Đặc tả\\\midrule\endfirsthead
\toprule Trường & Đặc tả\\\midrule\endhead
\textbf{ID / tác nhân / ưu tiên} & UC02.1; Chủ xe; Ưu tiên cao.\\
\textbf{Yêu cầu / kích hoạt} & FR02; Chọn Thêm xe.\\
\textbf{Tiền điều kiện} & Phiên hợp lệ; model/brand/type đang active.\\
\textbf{Luồng chính} & \begin{enumerate}[leftmargin=*,nosep]
\item Đọc loại xe, chọn hãng rồi dòng theo cascade.
\item Nhập tên tự đặt, plate, year, ngày baseline và ODO ban đầu; ImageUrl tùy chọn.
\item Server lấy owner từ claim, kiểm tra model active toàn chuỗi, chuẩn hóa plate.
\item Kiểm tra unique theo owner kể cả archived; ghi vehicle.
\item Trả 201 DTO với type/brand/model suy ra; refetch danh sách và chi tiết.
\end{enumerate}\\
\textbf{Luồng thay thế} & Đổi loại/hãng phải xóa lựa chọn con không tương thích. Có thể dùng ảnh mặc định khi không nhập URL.\\
\textbf{Ngoại lệ} & Thiếu trường/inactive catalog: 400; plate trùng: 409. Race unique phải được ánh xạ, không lộ SQL.\\
\textbf{Hậu điều kiện / chấp nhận} & Xe thuộc đúng owner và một model hợp lệ; chưa phát sinh chi phí.\\
\textbf{Kiểm thử liên quan} & AT02, AT03, AT04.\\
\bottomrule\end{longtable}\end{smalltable}
\subsubsection{UC02.2 -- Sửa thông tin và mốc xe}\label{uc:UC02.2}
\begin{smalltable}\begin{longtable}{L{3.2cm}L{11.8cm}}
\caption{UC02.2 -- Sửa thông tin và mốc xe}\label{tab:UC02.2}\\
\toprule Trường & Đặc tả\\\midrule\endfirsthead
\toprule Trường & Đặc tả\\\midrule\endhead
\textbf{ID / tác nhân / ưu tiên} & UC02.2; Chủ xe; Ưu tiên cao.\\
\textbf{Yêu cầu / kích hoạt} & FR02,FR03; Mở Sửa xe.\\
\textbf{Tiền điều kiện} & Xe active thuộc owner; có version mới nhất.\\
\textbf{Luồng chính} & \begin{enumerate}[leftmargin=*,nosep]
\item Tải thông tin và version; hiển thị catalog hiện tại.
\item Chỉnh tên/plate/model/year/ảnh hoặc baseline theo trường cho phép.
\item Server kiểm tra owner, active, version và model active.
\item Đổi baseline phải ValidateBaseOdo với lịch sử liên quan; kiểm tra plate không trùng.
\item Tăng version, lưu với concurrency token và trả DTO; UI cập nhật sau commit.
\end{enumerate}\\
\textbf{Luồng thay thế} & Hủy bỏ không ghi. Model cũ inactive vẫn cần hiển thị để hiểu lịch sử nhưng không chọn cho dữ liệu mới.\\
\textbf{Ngoại lệ} & Version cũ: 409; xe khác owner: 404; baseline vượt history bị từ chối. Form giữ giá trị để sửa.\\
\textbf{Hậu điều kiện / chấp nhận} & Thông tin mới nhất quán với chuỗi và catalog; lịch sử không bị gán chủ khác.\\
\textbf{Kiểm thử liên quan} & AT04, AT11.\\
\bottomrule\end{longtable}\end{smalltable}
\subsubsection{UC02.3 -- Lưu trữ và khôi phục phương tiện}\label{uc:UC02.3}
\begin{smalltable}\begin{longtable}{L{3.2cm}L{11.8cm}}
\caption{UC02.3 -- Lưu trữ và khôi phục phương tiện}\label{tab:UC02.3}\\
\toprule Trường & Đặc tả\\\midrule\endfirsthead
\toprule Trường & Đặc tả\\\midrule\endhead
\textbf{ID / tác nhân / ưu tiên} & UC02.3; Chủ xe; Ưu tiên cao.\\
\textbf{Yêu cầu / kích hoạt} & FR02,FR09; Chọn Lưu trữ hoặc Khôi phục.\\
\textbf{Tiền điều kiện} & Xe thuộc owner; biết version.\\
\textbf{Luồng chính} & \begin{enumerate}[leftmargin=*,nosep]
\item Hiển thị xác nhận rõ lịch sử được giữ và thao tác ghi sẽ bị chặn.
\item Gửi archive/restore với version.
\item Server kiểm tra owner/version rồi đặt hoặc xóa ArchivedAtUtc, tăng version.
\item thay trạng thái nhắc và đánh giá lại kế hoạch khi restore.
\item Refetch danh sách; xe archived chỉ xuất hiện khi chọn gồm lưu trữ.
\end{enumerate}\\
\textbf{Luồng thay thế} & Hủy xác nhận không thay đổi. Báo cáo chỉ gồm archived khi người dùng chọn.\\
\textbf{Ngoại lệ} & 404 cho sai owner; 409 cho version cũ. Mutation hoạt động trên xe archived phải bị chặn.\\
\textbf{Hậu điều kiện / chấp nhận} & Không xóa cứng xe hay lịch sử; restore không tạo completion giả.\\
\textbf{Kiểm thử liên quan} & AT03, AT09.\\
\bottomrule\end{longtable}\end{smalltable}
\subsubsection{UC03.1 -- Thêm ODO thủ công}\label{uc:UC03.1}
\begin{smalltable}\begin{longtable}{L{3.2cm}L{11.8cm}}
\caption{UC03.1 -- Thêm ODO thủ công}\label{tab:UC03.1}\\
\toprule Trường & Đặc tả\\\midrule\endfirsthead
\toprule Trường & Đặc tả\\\midrule\endhead
\textbf{ID / tác nhân / ưu tiên} & UC03.1; Chủ xe; Ưu tiên cao.\\
\textbf{Yêu cầu / kích hoạt} & FR03; Chọn Cập nhật ODO.\\
\textbf{Tiền điều kiện} & Xe active; ngày nghiệp vụ và ODO được biết.\\
\textbf{Luồng chính} & \begin{enumerate}[leftmargin=*,nosep]
\item Nhập ngày, ODO và ghi chú; UI ghi rõ km và ngày cập nhật.
\item Server đọc xe theo owner, tải history và baseline.
\item Kiểm tra ngày không tương lai, ODO nằm giữa reading trước/sau; cùng ngày cho bằng nhau.
\item Ghi Entry kind Odo với cost/liters bằng 0.
\item Refetch current ODO; thiết kế đích đánh giá lại trạng thái và nhắc.
\end{enumerate}\\
\textbf{Luồng thay thế} & Cho nhập quá khứ nếu chuỗi hợp lệ; reading cũ không thay current bằng giá trị thấp hơn.\\
\textbf{Ngoại lệ} & ODO âm/sai chuỗi/ngày tương lai: validation; archived bị chặn. Hai create đồng thời cần guard xe.\\
\textbf{Hậu điều kiện / chấp nhận} & Có một manual reading hợp lệ, không sinh chi phí.\\
\textbf{Kiểm thử liên quan} & AT04, AT11.\\
\bottomrule\end{longtable}\end{smalltable}
\subsubsection{UC03.2 -- Sửa hoặc xóa ODO thủ công}\label{uc:UC03.2}
\begin{smalltable}\begin{longtable}{L{3.2cm}L{11.8cm}}
\caption{UC03.2 -- Sửa hoặc xóa ODO thủ công}\label{tab:UC03.2}\\
\toprule Trường & Đặc tả\\\midrule\endfirsthead
\toprule Trường & Đặc tả\\\midrule\endhead
\textbf{ID / tác nhân / ưu tiên} & UC03.2; Chủ xe; Ưu tiên cao.\\
\textbf{Yêu cầu / kích hoạt} & FR03; Mở reading và chọn Sửa hoặc Xóa.\\
\textbf{Tiền điều kiện} & Reading manual của xe active; version hiện có.\\
\textbf{Luồng chính} & \begin{enumerate}[leftmargin=*,nosep]
\item Server tìm reading manual theo owner + vehicle + reading ID.
\item Nếu sửa, nhập ngày/ODO/note; loại reading hiện tại khỏi tập so sánh.
\item Validate chuỗi hoặc xác nhận xóa; kiểm tra version.
\item Lưu thay đổi với concurrency token; xóa reading thủ công không làm thay đổi hoạt động nguồn khác.
\item Tính current từ phần còn lại; thiết kế đích cập nhật hạn/nhắc.
\end{enumerate}\\
\textbf{Luồng thay thế} & Reading từ nhiên liệu/dịch vụ dẫn tới form nguồn; baseline xe dẫn tới Sửa xe.\\
\textbf{Ngoại lệ} & 409 nếu stale; 404 nếu không phải manual hoặc khác owner; sai chuỗi không được ghi.\\
\textbf{Hậu điều kiện / chấp nhận} & Không đọc mồ côi; xóa max có thể hạ current về mốc hợp lệ còn lại.\\
\textbf{Kiểm thử liên quan} & AT04, AT11.\\
\bottomrule\end{longtable}\end{smalltable}
\subsubsection{UC04.1 -- Ghi nhiên liệu}\label{uc:UC04.1}
\begin{smalltable}\begin{longtable}{L{3.2cm}L{11.8cm}}
\caption{UC04.1 -- Ghi nhiên liệu}\label{tab:UC04.1}\\
\toprule Trường & Đặc tả\\\midrule\endfirsthead
\toprule Trường & Đặc tả\\\midrule\endhead
\textbf{ID / tác nhân / ưu tiên} & UC04.1; Chủ xe; Ưu tiên cao.\\
\textbf{Yêu cầu / kích hoạt} & FR04; Chọn Thêm nhiên liệu.\\
\textbf{Tiền điều kiện} & Xe active; có ngày/ODO/lít/tiền; API đích được triển khai.\\
\textbf{Luồng chính} & \begin{enumerate}[leftmargin=*,nosep]
\item Nhập ngày, ODO, lít dương tối đa ba số lẻ và tiền VNĐ nguyên dương.
\item UI tạo RequestId và giữ ID đó cho cùng lần gửi/retry.
\item Server kiểm tra owner, active, số và chuỗi; tìm replay trước khi ghi.
\item Trong transaction có guard xe, ghi FuelRecord và linked reading; cập nhật dữ liệu liên quan.
\item Commit rồi trả tài nguyên; invalidate history, xe và report.
\end{enumerate}\\
\textbf{Luồng thay thế} & Cùng key/payload trả kết quả đã tạo; không tạo chi phí thứ hai.\\
\textbf{Ngoại lệ} & Khác payload cùng key: 409; lỗi giữa ghi rollback; mất response thì retry cùng ID.\\
\textbf{Hậu điều kiện / chấp nhận} & Một nguồn nhiên liệu và reading liên kết; tổng tiền tăng đúng một lần.\\
\textbf{Kiểm thử liên quan} & AT05, AT10.\\
\bottomrule\end{longtable}\end{smalltable}
\subsubsection{UC04.2 -- Điều chỉnh hoặc xóa nhiên liệu}\label{uc:UC04.2}
\begin{smalltable}\begin{longtable}{L{3.2cm}L{11.8cm}}
\caption{UC04.2 -- Điều chỉnh hoặc xóa nhiên liệu}\label{tab:UC04.2}\\
\toprule Trường & Đặc tả\\\midrule\endfirsthead
\toprule Trường & Đặc tả\\\midrule\endhead
\textbf{ID / tác nhân / ưu tiên} & UC04.2; Chủ xe; Ưu tiên cao.\\
\textbf{Yêu cầu / kích hoạt} & FR04,FR03; Mở lịch sử nhiên liệu.\\
\textbf{Tiền điều kiện} & Nguồn nhiên liệu thuộc xe active; version.\\
\textbf{Luồng chính} & \begin{enumerate}[leftmargin=*,nosep]
\item Tải nguồn và version; cho sửa trường allowlist hoặc xác nhận xóa.
\item Kiểm tra owner, active, version và reading nguồn.
\item Guard xe; validate chuỗi sau thay đổi.
\item Trong cùng transaction sửa/xóa fuel và reading, đánh giá lại plan bị ảnh hưởng ODO.
\item Commit rồi làm mới lịch sử và báo cáo.
\end{enumerate}\\
\textbf{Luồng thay thế} & Hủy không ghi. Khi xóa fuel không xóa dịch vụ hoặc kế hoạch cùng ngày.\\
\textbf{Ngoại lệ} & Stale 409; invalid ODO không commit; lỗi DB giữ nguyên nguồn/reading.\\
\textbf{Hậu điều kiện / chấp nhận} & Tổng báo cáo phản ánh dữ liệu còn lại; không có reading mồ côi.\\
\textbf{Kiểm thử liên quan} & AT05, AT11.\\
\bottomrule\end{longtable}\end{smalltable}
\subsubsection{UC05.1 -- Ghi dịch vụ nhiều dòng và completion}\label{uc:UC05.1}
\begin{smalltable}\begin{longtable}{L{3.2cm}L{11.8cm}}
\caption{UC05.1 -- Ghi dịch vụ nhiều dòng và completion}\label{tab:UC05.1}\\
\toprule Trường & Đặc tả\\\midrule\endfirsthead
\toprule Trường & Đặc tả\\\midrule\endhead
\textbf{ID / tác nhân / ưu tiên} & UC05.1; Chủ xe; Ưu tiên cao.\\
\textbf{Yêu cầu / kích hoạt} & FR05,FR08; Thêm dịch vụ hoặc Ghi đã thực hiện từ nhắc.\\
\textbf{Tiền điều kiện} & Xe active; plan liên kết nếu có cùng xe/owner.\\
\textbf{Luồng chính} & \begin{enumerate}[leftmargin=*,nosep]
\item Nhập ngày/ODO và ít nhất một item gồm loại, tên và tiền.
\item Item Maintenance tùy chọn plan; Repair không có plan; không lặp plan trong cùng dịch vụ.
\item Server kiểm tra tổng dương, baseline, uniqueness completion cùng ngày và RequestId.
\item Guard xe và transaction: ghi service/items/reading; tìm latest completion mỗi plan bị ảnh hưởng.
\item Tính kỳ và nhắc; commit; hiển thị chi phí và hạn mới.
\end{enumerate}\\
\textbf{Luồng thay thế} & Mở từ nhắc điền sẵn xe/plan nhưng cho kiểm tra trước lưu. Dịch vụ không liên kết vẫn lưu lịch sử.\\
\textbf{Ngoại lệ} & Plan sai xe/trước baseline/trùng ngày: từ chối; lỗi một item rollback mọi phần.\\
\textbf{Hậu điều kiện / chấp nhận} & Chỉ plan được liên kết chuyển kỳ; tiền dịch vụ tính một lần.\\
\textbf{Kiểm thử liên quan} & AT06, AT08, AT10.\\
\bottomrule\end{longtable}\end{smalltable}
\subsubsection{UC05.2 -- Sửa/xóa dịch vụ và tính lại kỳ}\label{uc:UC05.2}
\begin{smalltable}\begin{longtable}{L{3.2cm}L{11.8cm}}
\caption{UC05.2 -- Sửa/xóa dịch vụ và tính lại kỳ}\label{tab:UC05.2}\\
\toprule Trường & Đặc tả\\\midrule\endfirsthead
\toprule Trường & Đặc tả\\\midrule\endhead
\textbf{ID / tác nhân / ưu tiên} & UC05.2; Chủ xe; Ưu tiên cao.\\
\textbf{Yêu cầu / kích hoạt} & FR05,FR08; Chọn Sửa/Xóa từ lịch sử nguồn.\\
\textbf{Tiền điều kiện} & Dịch vụ xe active; có version và các plan liên quan.\\
\textbf{Luồng chính} & \begin{enumerate}[leftmargin=*,nosep]
\item Tải service/items và version; giữ tập plan trước sửa.
\item Nhập thay đổi hoặc xác nhận xóa; kiểm tra quyền và chuỗi.
\item Guard xe, cập nhật nguồn/items/reading trong transaction.
\item Lấy hợp của plan trước và sau sửa; tìm latest completion còn hợp lệ cho từng plan.
\item Tính lại due/notification; commit rồi làm mới UI.
\end{enumerate}\\
\textbf{Luồng thay thế} & Bỏ liên kết không xóa plan. Khi không còn completion dùng baseline.\\
\textbf{Ngoại lệ} & Plan liên kết không hợp lệ hoặc stale 409: không ghi phần nào. Không chọn mốc theo thời điểm request.\\
\textbf{Hậu điều kiện / chấp nhận} & Kỳ dựa trên lịch sử còn lại; tổng tiền giảm/đổi đúng số liệu nguồn.\\
\textbf{Kiểm thử liên quan} & AT08, AT11.\\
\bottomrule\end{longtable}\end{smalltable}
\subsubsection{UC06.1 -- Tạo hoặc sửa kế hoạch bảo dưỡng}\label{uc:UC06.1}
\begin{smalltable}\begin{longtable}{L{3.2cm}L{11.8cm}}
\caption{UC06.1 -- Tạo hoặc sửa kế hoạch bảo dưỡng}\label{tab:UC06.1}\\
\toprule Trường & Đặc tả\\\midrule\endfirsthead
\toprule Trường & Đặc tả\\\midrule\endhead
\textbf{ID / tác nhân / ưu tiên} & UC06.1; Chủ xe; Ưu tiên cao.\\
\textbf{Yêu cầu / kích hoạt} & FR06,FR07; Chọn Thêm/Sửa kế hoạch.\\
\textbf{Tiền điều kiện} & Xe active; có baseline từ dữ liệu đáng tin cậy.\\
\textbf{Luồng chính} & \begin{enumerate}[leftmargin=*,nosep]
\item Nhập tên, mode Time/Distance/Either và ý nghĩa baseline.
\item Hiện trường đúng mode: ngày/tháng/ngưỡng ngày hoặc ODO/km/ngưỡng km.
\item Server kiểm tra baseline, chu kỳ, ngưỡng, tên unique trong xe và completion liên quan.
\item Lưu trong transaction, đánh giá trạng thái từ dữ liệu hiện hành.
\item Trả due, remaining, reason và cờ thiếu/cũ; UI giải thích điều kiện đến trước.
\end{enumerate}\\
\textbf{Luồng thay thế} & Baseline là bắt đầu theo dõi không sinh hóa đơn; sửa baseline khi có completion phải không phá lịch sử.\\
\textbf{Ngoại lệ} & Chu kỳ 0, ngưỡng vượt kỳ, ngày tương lai hoặc ODO vượt current bị từ chối.\\
\textbf{Hậu điều kiện / chấp nhận} & Kế hoạch đủ cấu hình cho mode; không tạo chi phí giả.\\
\textbf{Kiểm thử liên quan} & AT07.\\
\bottomrule\end{longtable}\end{smalltable}
\subsubsection{UC06.2 -- Tạm dừng, tiếp tục hoặc xóa kế hoạch}\label{uc:UC06.2}
\begin{smalltable}\begin{longtable}{L{3.2cm}L{11.8cm}}
\caption{UC06.2 -- Tạm dừng, tiếp tục hoặc xóa kế hoạch}\label{tab:UC06.2}\\
\toprule Trường & Đặc tả\\\midrule\endfirsthead
\toprule Trường & Đặc tả\\\midrule\endhead
\textbf{ID / tác nhân / ưu tiên} & UC06.2; Chủ xe; Ưu tiên cao.\\
\textbf{Yêu cầu / kích hoạt} & FR06,FR09; Chọn Pause/Resume/Xóa.\\
\textbf{Tiền điều kiện} & Plan thuộc xe active; version; biết đã có completion hay chưa.\\
\textbf{Luồng chính} & \begin{enumerate}[leftmargin=*,nosep]
\item Hiển thị tác động: pause dừng nhắc chủ động, không xóa lịch sử.
\item Server kiểm tra owner/version và trạng thái xe.
\item Pause/resume giữ baseline và lịch sử; resume đánh giá dữ liệu hiện hành.
\item Chỉ cho xóa plan chưa có completion; đã có lịch sử hướng dẫn pause.
\item Commit và cập nhật danh sách/badge.
\end{enumerate}\\
\textbf{Luồng thay thế} & Hủy xác nhận không đổi. Plan paused vẫn tra cứu được với nhãn riêng.\\
\textbf{Ngoại lệ} & Xóa plan đã có completion bị từ chối; stale 409; sai owner 404.\\
\textbf{Hậu điều kiện / chấp nhận} & Không reset kỳ; không làm mất dịch vụ/chi phí.\\
\textbf{Kiểm thử liên quan} & AT07, AT09.\\
\bottomrule\end{longtable}\end{smalltable}
\subsubsection{UC07.1 -- Xem hạn và giải thích trạng thái}\label{uc:UC07.1}
\begin{smalltable}\begin{longtable}{L{3.2cm}L{11.8cm}}
\caption{UC07.1 -- Xem hạn và giải thích trạng thái}\label{tab:UC07.1}\\
\toprule Trường & Đặc tả\\\midrule\endfirsthead
\toprule Trường & Đặc tả\\\midrule\endhead
\textbf{ID / tác nhân / ưu tiên} & UC07.1; Chủ xe; Ưu tiên cao.\\
\textbf{Yêu cầu / kích hoạt} & FR07; Mở kế hoạch hoặc nhắc.\\
\textbf{Tiền điều kiện} & Có quyền xe/plan; query không cần xe active để xem.\\
\textbf{Luồng chính} & \begin{enumerate}[leftmargin=*,nosep]
\item Đọc mode, mốc latest completion hoặc baseline và current ODO.
\item Evaluator tính hạn tháng lịch và/hoặc ODO.
\item So từng chiều với todayVN/current; lấy mức ưu tiên đúng cho Either.
\item Trả trạng thái có chữ, lý do, remaining và missing-data flags.
\item UI hiện thêm paused/archive/ODO updatedAt/stale theo dữ liệu riêng.
\end{enumerate}\\
\textbf{Luồng thay thế} & Thiếu một chiều nhưng chiều khác đã Soon/Due/Overdue vẫn cảnh báo; chỉ NotDue thì Unknown.\\
\textbf{Ngoại lệ} & Không đủ dữ liệu không được hiện mọi thứ ổn. ODO cũ không là chứng cứ chưa vượt hạn.\\
\textbf{Hậu điều kiện / chấp nhận} & Người dùng giải thích được nguồn hạn; không thay đổi completion.\\
\textbf{Kiểm thử liên quan} & AT07.\\
\bottomrule\end{longtable}\end{smalltable}
\subsubsection{UC07.2 -- Đọc inbox và đánh dấu đã đọc}\label{uc:UC07.2}
\begin{smalltable}\begin{longtable}{L{3.2cm}L{11.8cm}}
\caption{UC07.2 -- Đọc inbox và đánh dấu đã đọc}\label{tab:UC07.2}\\
\toprule Trường & Đặc tả\\\midrule\endfirsthead
\toprule Trường & Đặc tả\\\midrule\endhead
\textbf{ID / tác nhân / ưu tiên} & UC07.2; Chủ xe; Ưu tiên cao.\\
\textbf{Yêu cầu / kích hoạt} & FR09; Mở biểu tượng nhắc hoặc một thông báo.\\
\textbf{Tiền điều kiện} & Có notification của owner trong DB ở thiết kế đích.\\
\textbf{Luồng chính} & \begin{enumerate}[leftmargin=*,nosep]
\item Lấy inbox theo owner và trạng thái Current/Superseded, có phân trang.
\item Hiển thị xe, plan, level, lý do và thời điểm tạo.
\item Khi chọn đã đọc, cập nhật ReadAt idempotent.
\item Làm mới badge chỉ đếm Current và Unread.
\item Nếu đã làm bảo dưỡng, chuyển UC05.1 để ghi dịch vụ thật.
\end{enumerate}\\
\textbf{Luồng thay thế} & Từ chối browser permission vẫn dùng inbox; không buộc cấp quyền để xem dữ liệu.\\
\textbf{Ngoại lệ} & 404 cho nhắc khác owner; lỗi mạng giữ unread thực tế; quét worker không tạo bản trùng.\\
\textbf{Hậu điều kiện / chấp nhận} & Đọc chỉ đổi trạng thái đọc, không đổi kỳ hoặc chi phí.\\
\textbf{Kiểm thử liên quan} & AT09.\\
\bottomrule\end{longtable}\end{smalltable}
\subsubsection{UC08.1 -- Tra cứu lịch sử thống nhất}\label{uc:UC08.1}
\begin{smalltable}\begin{longtable}{L{3.2cm}L{11.8cm}}
\caption{UC08.1 -- Tra cứu lịch sử thống nhất}\label{tab:UC08.1}\\
\toprule Trường & Đặc tả\\\midrule\endfirsthead
\toprule Trường & Đặc tả\\\midrule\endhead
\textbf{ID / tác nhân / ưu tiên} & UC08.1; Chủ xe; Ưu tiên cao.\\
\textbf{Yêu cầu / kích hoạt} & FR10; Mở Lịch sử hoặc lịch sử của xe.\\
\textbf{Tiền điều kiện} & Có session; dữ liệu có thể rỗng.\\
\textbf{Luồng chính} & \begin{enumerate}[leftmargin=*,nosep]
\item Chọn xe, from/to, loại hoạt động, keyword và trang.
\item Server validate khoảng và owner; archived theo filter.
\item Query sort ổn định theo ngày và ID; reading sinh từ nhiên liệu hoặc dịch vụ không tạo thêm dòng hoạt động trùng.
\item Hiện dòng ngày/xe/loại/nội dung/ODO/tiền; mở chi tiết nguồn.
\item Chuyển sửa/xóa chỉ khi xe active và user có quyền.
\end{enumerate}\\
\textbf{Luồng thay thế} & Không kết quả: empty state kèm xóa lọc. Chuyển trang giữ bộ lọc.\\
\textbf{Ngoại lệ} & from > to hoặc page sai: validation; lỗi mạng có retry; không rò số lượng tài khoản khác.\\
\textbf{Hậu điều kiện / chấp nhận} & Query không ghi nghiệp vụ; một fuel/service chỉ là một hoạt động.\\
\textbf{Kiểm thử liên quan} & AT12, AT14.\\
\bottomrule\end{longtable}\end{smalltable}
\subsubsection{UC09.1 -- Xem dashboard}\label{uc:UC09.1}
\begin{smalltable}\begin{longtable}{L{3.2cm}L{11.8cm}}
\caption{UC09.1 -- Xem dashboard}\label{tab:UC09.1}\\
\toprule Trường & Đặc tả\\\midrule\endfirsthead
\toprule Trường & Đặc tả\\\midrule\endhead
\textbf{ID / tác nhân / ưu tiên} & UC09.1; Chủ xe; Ưu tiên cao.\\
\textbf{Yêu cầu / kích hoạt} & FR11; Mở trang Tổng quan.\\
\textbf{Tiền điều kiện} & Phiên hợp lệ; dữ liệu cùng owner.\\
\textbf{Luồng chính} & \begin{enumerate}[leftmargin=*,nosep]
\item Chọn một/tất cả xe và khoảng/tháng phù hợp.
\item Query dữ liệu tổng hợp dùng cùng điều kiện owner/date/archive.
\item Hiện số xe, chi phí, trạng thái cần chú ý và hoạt động gần đây.
\item Từ thẻ chỉ số, đi đến lịch sử hoặc report với filter được giữ.
\item Sau mutation liên quan, invalidate query để không hiển thị tổng cũ.
\end{enumerate}\\
\textbf{Luồng thay thế} & Tài khoản mới có zero/empty và CTA thêm xe; dữ liệu mẫu được tách khỏi số liệu cá nhân.\\
\textbf{Ngoại lệ} & Loading/error phân biệt zero thực tế; không coi lỗi fetch là mọi thứ ổn.\\
\textbf{Hậu điều kiện / chấp nhận} & Chỉ số đối chiếu được với nguồn; không phát sinh bản ghi.\\
\textbf{Kiểm thử liên quan} & AT10, AT12.\\
\bottomrule\end{longtable}\end{smalltable}
\subsubsection{UC09.2 -- Lập báo cáo chi phí}\label{uc:UC09.2}
\begin{smalltable}\begin{longtable}{L{3.2cm}L{11.8cm}}
\caption{UC09.2 -- Lập báo cáo chi phí}\label{tab:UC09.2}\\
\toprule Trường & Đặc tả\\\midrule\endfirsthead
\toprule Trường & Đặc tả\\\midrule\endhead
\textbf{ID / tác nhân / ưu tiên} & UC09.2; Chủ xe; Ưu tiên cao.\\
\textbf{Yêu cầu / kích hoạt} & FR11; Mở Báo cáo.\\
\textbf{Tiền điều kiện} & Dữ liệu có thể rỗng; cần khoảng ngày hợp lệ.\\
\textbf{Luồng chính} & \begin{enumerate}[leftmargin=*,nosep]
\item Chọn xe, ngày hai đầu, nhóm chi phí và tùy chọn archived.
\item Server kiểm tra owner; không nhận owner tùy ý từ browser.
\item Tổng hợp fuel amount và service item amount, nhóm theo tháng/xe/loại.
\item Trả tổng, series và rows cùng filter; UI trình bày bảng kèm biểu đồ.
\item Đối chiếu fixture 200.000 + 100.000 + 150.000 = 450.000.
\end{enumerate}\\
\textbf{Luồng thay thế} & Không dữ liệu trả 0 và empty series hợp lệ; mặc định không gồm archived.\\
\textbf{Ngoại lệ} & Khoảng đảo bị từ chối; tiền không dùng floating point để tích lũy; không cộng header service lần hai.\\
\textbf{Hậu điều kiện / chấp nhận} & Dashboard/bảng/biểu đồ nhất quán sau refresh; các khoản tiền đều có thể đối chiếu với bản ghi nguồn.\\
\textbf{Kiểm thử liên quan} & AT10, AT13.\\
\bottomrule\end{longtable}\end{smalltable}
\subsubsection{UC10.1 -- Xem và sửa hồ sơ}\label{uc:UC10.1}
\begin{smalltable}\begin{longtable}{L{3.2cm}L{11.8cm}}
\caption{UC10.1 -- Xem và sửa hồ sơ}\label{tab:UC10.1}\\
\toprule Trường & Đặc tả\\\midrule\endfirsthead
\toprule Trường & Đặc tả\\\midrule\endhead
\textbf{ID / tác nhân / ưu tiên} & UC10.1; Chủ xe; Ưu tiên cao.\\
\textbf{Yêu cầu / kích hoạt} & FR01; Mở Hồ sơ.\\
\textbf{Tiền điều kiện} & Có phiên; email chỉ đọc.\\
\textbf{Luồng chính} & \begin{enumerate}[leftmargin=*,nosep]
\item GET hồ sơ account hiện tại từ session.
\item Chỉnh tên và số điện thoại; email hiển thị ở chế độ chỉ đọc.
\item Server allowlist field, kiểm tra tên/phone và version nếu contract có.
\item Lưu thông tin cùng owner và trả DTO an toàn.
\item UI cập nhật tên hiển thị sau commit.
\end{enumerate}\\
\textbf{Luồng thay thế} & Hủy bỏ thay đổi; email không có nút đổi trong v1.\\
\textbf{Ngoại lệ} & Tên rỗng/phone sai bị từ chối; payload email/role/accountId không được mass assign.\\
\textbf{Hậu điều kiện / chấp nhận} & Chỉ profile được sửa; quyền/identity không đổi.\\
\textbf{Kiểm thử liên quan} & AT01, AT02.\\
\bottomrule\end{longtable}\end{smalltable}
\subsubsection{UC10.2 -- Đổi mật khẩu}\label{uc:UC10.2}
\begin{smalltable}\begin{longtable}{L{3.2cm}L{11.8cm}}
\caption{UC10.2 -- Đổi mật khẩu}\label{tab:UC10.2}\\
\toprule Trường & Đặc tả\\\midrule\endfirsthead
\toprule Trường & Đặc tả\\\midrule\endhead
\textbf{ID / tác nhân / ưu tiên} & UC10.2; Chủ xe; Ưu tiên cao.\\
\textbf{Yêu cầu / kích hoạt} & FR01,NFR02; Chọn Đổi mật khẩu.\\
\textbf{Tiền điều kiện} & Có phiên; biết password cũ.\\
\textbf{Luồng chính} & \begin{enumerate}[leftmargin=*,nosep]
\item Nhập password hiện tại, password mới và xác nhận.
\item Server kiểm tra CSRF và password cũ qua Identity.
\item Áp policy password mới; cập nhật security state trong transaction phù hợp.
\item Thu hồi các phiên theo policy đích: tất cả phiên của tài khoản, gồm phiên hiện tại.
\item Trả thành công rồi xóa cache và yêu cầu đăng nhập lại.
\end{enumerate}\\
\textbf{Luồng thay thế} & Hủy không đổi; không cung cấp reset qua email trong v1.\\
\textbf{Ngoại lệ} & Password cũ sai hoặc policy không đạt không đổi hash/session. Không log bất kỳ password nào.\\
\textbf{Hậu điều kiện / chấp nhận} & Password cũ không login được, phiên cũ bị từ chối.\\
\textbf{Kiểm thử liên quan} & AT01, AT16.\\
\bottomrule\end{longtable}\end{smalltable}
\subsubsection{UC10.3 -- Điều chỉnh nhắc và ngưỡng ODO cũ}\label{uc:UC10.3}
\begin{smalltable}\begin{longtable}{L{3.2cm}L{11.8cm}}
\caption{UC10.3 -- Điều chỉnh nhắc và ngưỡng ODO cũ}\label{tab:UC10.3}\\
\toprule Trường & Đặc tả\\\midrule\endfirsthead
\toprule Trường & Đặc tả\\\midrule\endhead
\textbf{ID / tác nhân / ưu tiên} & UC10.3; Chủ xe; Ưu tiên cao.\\
\textbf{Yêu cầu / kích hoạt} & FR12,FR09; Mở Cài đặt.\\
\textbf{Tiền điều kiện} & Có session; settings của owner.\\
\textbf{Luồng chính} & \begin{enumerate}[leftmargin=*,nosep]
\item Đọc remindersEnabled và staleOdoDays cùng cấu hình được hỗ trợ.
\item Chọn bật/tắt, nhập staleOdoDays nguyên dương.
\item Server validate/allowlist; lưu đúng owner.
\item Tắt nhắc không xóa plan/lịch sử; loại nhắc không còn current theo policy.
\item Bật lại đánh giá trạng thái hiện hành, không tạo nhắc cho mọi kỳ quá khứ.
\end{enumerate}\\
\textbf{Luồng thay thế} & Browser permission được hướng dẫn riêng, không đồng nghĩa cấu hình nhắc server.\\
\textbf{Ngoại lệ} & Số âm/0/phần lẻ bị từ chối. Lỗi lưu không đổi giá trị thật bằng UI optimistic sai.\\
\textbf{Hậu điều kiện / chấp nhận} & Cấu hình bền vững, ảnh hưởng đúng lần đánh giá tiếp theo.\\
\textbf{Kiểm thử liên quan} & AT15.\\
\bottomrule\end{longtable}\end{smalltable}
\subsubsection{UC11.1 -- Đọc danh mục loại--hãng--dòng}\label{uc:UC11.1}
\begin{smalltable}\begin{longtable}{L{3.2cm}L{11.8cm}}
\caption{UC11.1 -- Đọc danh mục loại--hãng--dòng}\label{tab:UC11.1}\\
\toprule Trường & Đặc tả\\\midrule\endfirsthead
\toprule Trường & Đặc tả\\\midrule\endhead
\textbf{ID / tác nhân / ưu tiên} & UC11.1; Khách / chủ xe; hỗ trợ FR02; ưu tiên cao.\\
\textbf{Yêu cầu / kích hoạt} & FR02 mở rộng; Mở form xe.\\
\textbf{Tiền điều kiện} & Có quyền đọc catalog theo controller; catalog có thể rỗng.\\
\textbf{Luồng chính} & \begin{enumerate}[leftmargin=*,nosep]
\item GET types active.
\item Chọn type, GET brands thuộc type; xóa brand/model cũ khi type đổi.
\item Chọn brand, GET models thuộc brand.
\item Chọn modelId và gửi cùng vehicle form.
\item Server suy ra type/brand từ model, không tin các tên do UI gửi.
\end{enumerate}\\
\textbf{Luồng thay thế} & Không có model: empty state, không cho gửi giá trị tự do thay ID.\\
\textbf{Ngoại lệ} & ID không tồn tại/inactive bị từ chối lúc ghi; không gộp model khác brand.\\
\textbf{Hậu điều kiện / chấp nhận} & Vehicle tham chiếu một model; quan hệ đúng được bảo đảm server và FK.\\
\textbf{Kiểm thử liên quan} & AT03, AT16.\\
\bottomrule\end{longtable}\end{smalltable}
\subsubsection{UC11.2 -- Quản trị catalog}\label{uc:UC11.2}
\begin{smalltable}\begin{longtable}{L{3.2cm}L{11.8cm}}
\caption{UC11.2 -- Quản trị catalog}\label{tab:UC11.2}\\
\toprule Trường & Đặc tả\\\midrule\endfirsthead
\toprule Trường & Đặc tả\\\midrule\endhead
\textbf{ID / tác nhân / ưu tiên} & UC11.2; Admin catalog; quản lý danh mục; ưu tiên cao.\\
\textbf{Yêu cầu / kích hoạt} & FR02 mở rộng; Gọi API tạo/cập nhật metadata.\\
\textbf{Tiền điều kiện} & Account đã được cấp role Admin hợp lệ; không có seed Admin được xác nhận.\\
\textbf{Luồng chính} & \begin{enumerate}[leftmargin=*,nosep]
\item Xác thực và kiểm tra role Admin ở controller.
\item Tạo/sửa type code/name hoặc brand/model đúng parent.
\item Kiểm tra tên unique trong parent, version và quan hệ.
\item Khi ngừng dùng đặt IsActive=false; không phá vehicle lịch sử.
\item Trả DTO và cập nhật danh mục cho thao tác sau.
\end{enumerate}\\
\textbf{Luồng thay thế} & Không làm admin dashboard/user management trong v1; chức năng này có thể thao tác qua API khi được cấp quyền đúng.\\
\textbf{Ngoại lệ} & Không role: 403; trùng tên/version cũ: conflict; parent không phù hợp: validation.\\
\textbf{Hậu điều kiện / chấp nhận} & Metadata đúng quan hệ; user thường không thay danh mục chung.\\
\textbf{Kiểm thử liên quan} & AT03, AT16.\\
\bottomrule\end{longtable}\end{smalltable}


\section{Quy tắc nghiệp vụ}\label{sec:business-rules}
\subsection{BR01 -- Quyền sở hữu}
Account là từ claim đã xác thực. Mọi lookup theo tài nguyên riêng có điều kiện owner; activity, reading, plan và từng item liên kết phải cùng xe. ID tồn tại ở tài khoản khác trả 404 như không tồn tại, không trả tên xe hoặc metadata. Catalog dùng chung có scope khác; Admin catalog không được bỏ kiểm tra owner ở API xe.
\subsection{BR02 -- Ngày và chuỗi ODO}
Ngày nghiệp vụ không sau ngày hiện tại Việt Nam, độc lập timezone OS. Timestamp audit dùng UTC. Số km được ghi bằng số nguyên không âm; chuỗi tăng không giảm giữa các ngày; cùng ngày cho bằng nhau và sắp ODO tăng, thêm ID để ổn định. Khi chèn/sửa, loại chính reading đang sửa rồi kiểm tra max ngày trước và min ngày sau. Current ODO là ODO lớn nhất ở ngày mới nhất; với chuỗi hợp lệ không giảm, cũng bằng max toàn chuỗi và baseline.

Baseline xe không xóa bằng endpoint manual; đổi qua sửa xe và kiểm tra toàn chuỗi. Xóa reading max làm current ODO giảm về reading còn lại nếu hợp lệ. Fuel/service reading chỉ sửa/xóa qua nguồn.
\subsection{BR03 -- Baseline và cấu hình plan}
Time cần ngày baseline, chu kỳ tháng nguyên dương và ngưỡng ngày không âm. Distance cần baseline ODO, chu kỳ km nguyên dương và ngưỡng km không âm. Either cần cả hai. Baseline phân biệt ``lần bảo dưỡng đã biết'' với ``bắt đầu theo dõi''; không tạo chi phí giả. Ngày không tương lai; ODO không vượt current và phải phù hợp lịch sử.

Ngưỡng km không vượt chu kỳ km. Ngưỡng ngày không vượt độ dài kỳ tương ứng tính theo tháng lịch; nếu đổi baseline/chu kỳ khiến ngưỡng không phù hợp phải báo lỗi, không silently clamp. Đối với chiều không dùng, field phải null/absent theo kiểu kế hoạch, tránh giữ giá trị ẩn ảnh hưởng evaluator.
\subsection{BR04 -- Hạn và mốc hoàn thành}
Chọn completion hợp lệ mới nhất theo ngày của plan; chưa có thì dùng baseline. Hạn ngày = ngày mốc cộng số tháng lịch, ngày cuối tháng chặn về ngày hợp lệ. Hạn ODO = ODO mốc + chu kỳ km. Không liên kết completion trước baseline ở chiều cấu hình. Dịch vụ cũ vẫn có thể lưu không liên kết. Không có hai completion của một plan cùng ngày; ràng buộc này được kiểm tra khi lưu dữ liệu.
\subsection{BR05 -- Trạng thái có thể giải thích}
Từng chiều dùng thứ tự: vượt hạn là Overdue; bằng hạn là Due; chưa đến nhưng còn trong ngưỡng là Soon; xa hơn là NotDue; thiếu dữ liệu là Unknown. Either ưu tiên Overdue $>$ Due $>$ Soon $>$ NotDue. Nếu thiếu một chiều, chiều còn lại Soon/Due/Overdue đủ để kết luận mức cần chú ý; nếu chỉ NotDue thì kết quả Unknown kèm lý do thiếu. Pause, archive và ODO cũ là các cờ riêng. Không coi ODO cũ là chứng cứ xe vẫn chưa đến hạn.
\subsection{BR06 -- Tính lại sau sửa/xóa}
Chỉ Maintenance có plan mới chuyển kỳ. Repair hoặc Maintenance không liên kết chỉ đóng góp lịch sử/chi phí. Sửa/xóa dịch vụ tìm lại latest completion của toàn bộ plan bị ảnh hưởng, gồm cả plan bị bỏ liên kết. Không lấy bản vừa sửa làm mốc chỉ vì nó được gửi gần nhất. Thay nguồn, reading, kỳ và notification phải nhất quán trong transaction; lỗi rollback.
\subsection{BR07 -- Tiền và báo cáo}
Tiền được ghi bằng VNĐ nguyên dưới dạng decimal/numeric; fuel amount dương, liters dương tối đa ba số lẻ. Dịch vụ có ít nhất một item, mỗi tiền không âm và tổng dương. Tổng bằng fuel amount cộng từng service item amount; không cộng thêm header total. Khoảng ngày bao gồm hai đầu; from $>$ to bị từ chối. Mặc định loại archive, có tùy chọn đưa vào. Biểu đồ, bảng và số tổng lấy cùng tập dữ liệu.
\subsection{BR08 -- Vòng đời nhắc}
Tối đa một notification cho plan/kỳ/mức. Kỳ định danh bằng mốc baseline/completion ổn định, không dùng tên plan hoặc thời điểm quét. Đổi tên/đọc không tạo kỳ. Nhắc cũ Superseded không tính badge. Pause/archive/tắt nhắc dừng tạo nhắc chủ động; resume/restore/bật lại đánh giá hiện hành. Đọc không hoàn thành plan.

Khi cùng một thông báo được kích hoạt lại, hệ thống giữ trạng thái đã đọc. Nếu kế hoạch chuyển sang mức cần chú ý cao hơn, hệ thống tạo thông báo cho mức mới và đánh dấu chưa đọc. Thông báo trình duyệt chỉ được phát khi có thay đổi phù hợp và người dùng đã cấp quyền.
\subsection{BR09 -- Retry và cạnh tranh}
Fuel/service create có RequestId theo owner; cùng ID và canonical payload trả tài nguyên đã tạo, payload khác trả 409. Kiểm tra replay trước điều kiện version để retry thành công không thành stale error. Khóa yêu cầu được lưu cùng kết quả của giao dịch để hỗ trợ việc thử lại sau khi mất kết nối. Thời gian lưu khóa được quy định trong hợp đồng của từng chức năng.

Update/delete có version. Mutation ảnh hưởng chuỗi và kỳ cần guard cấp xe trong transaction bằng concurrency token của xe hoặc cơ chế khóa phù hợp. Version từng reading chỉ chống ghi đè cùng reading, chưa ngăn hai create khác ID tạo chuỗi sai. Mức isolation/locking phải được kiểm thử bằng hai DbContext/connection thật.

\section{Validation và hợp đồng dữ liệu đầu vào}
Các trường dữ liệu được kiểm tra ở máy chủ trước khi lưu. Biểu mẫu hiển thị lỗi tại trường liên quan và giữ nội dung đã nhập để người dùng sửa lại.
\begin{smalltable}\begin{longtable}{L{2.1cm}L{8.6cm}L{3.8cm}}
\caption{Ma trận validation và phản hồi người dùng}\label{tab:validation}\\\toprule Đối tượng & Ràng buộc kiểm tra tại server & Lỗi / phục hồi\\\midrule\endfirsthead
\toprule Đối tượng & Ràng buộc kiểm tra tại server & Lỗi / phục hồi\\\midrule\endhead
Đăng ký & Email hợp lệ/unique; tên 2--60 ký tự; phone rỗng hoặc 10 số bắt đầu 0; password theo chính sách mật khẩu. & Lỗi theo trường, không log password; đăng ký xong đăng nhập riêng.\\
Xe & Tên 2--60; biển số 2--20 và chuẩn hóa unique theo owner; năm 1980--2100 và danh mục loại--hãng--dòng xe còn hoạt động; ngày và số km ban đầu hợp lệ. & Giữ input; đổi type/brand phải xóa lựa chọn con không còn hợp lệ.\\
Reading & Ngày không tương lai theo Việt Nam; ODO nguyên 0--9.999.999; phù hợp reading trước/sau; baseline và nguồn được bảo vệ. & Giải thích khoảng hợp lệ; đọc nguồn để sửa fuel/service reading.\\
Nhiên liệu & Title 2--100; note tối đa 500; lít dương, tối đa 500, không quá ba số lẻ; VNĐ nguyên dương tối đa 10 tỷ; RequestId khi tạo. & Một payload sai không tạo nguồn/reading; không tự suy tiền từ giá thị trường.\\
Dịch vụ & Ít nhất một item; mỗi item loại hợp lệ, có tên, tiền nguyên không âm tối đa 10 tỷ; tổng dương. Chỉ Maintenance link plan cùng xe. & Chỉ đúng dòng lỗi nhưng rollback toàn yêu cầu; không báo đã lưu một phần.\\
Kế hoạch & Mode và meaning đóng; trường đủ theo mode; chu kỳ dương; ngưỡng không âm/không vượt kỳ; baseline tương thích lịch sử; tên unique trong xe. & Chỉ rõ field lỗi; không để giá trị ẩn của mode trước ảnh hưởng evaluator.\\
Cài đặt & remindersEnabled là boolean; staleOdoDays nguyên dương, mặc định 30 ngày; kiểm tra phiên bản khi cập nhật. & Giữ giá trị thật khi save lỗi; phân biệt ODO cũ với ngưỡng sắp hạn.\\
Query & from không sau to; IDs thuộc owner; loại/nhóm/page/pageSize hợp lệ. Mỗi trang mặc định 20 bản ghi, tối đa 100 bản ghi. & Giữ filter khi empty/error; lỗi tải không thành tổng 0.\\\bottomrule
\end{longtable}\end{smalltable}
Ngày nghiệp vụ dùng \code{YYYY-MM-DD}; timestamp kỹ thuật UTC; JSON camelCase. Không nhận owner từ payload làm quyền truy cập. Số lẻ vượt precision bị từ chối, không âm thầm làm tròn.

\section{Yêu cầu phi chức năng}
Bên cạnh chức năng nghiệp vụ, nhóm xác định các yêu cầu về bảo mật, tính toàn vẹn, khả năng sử dụng và vận hành. Các tiêu chí sau là cơ sở xây dựng bài kiểm tra nghiệm thu.
\begin{smalltable}\begin{longtable}{L{1.3cm}L{3cm}L{6.2cm}L{4cm}}
\caption{NFR và phương pháp chứng minh}\label{tab:nfr}\\\toprule Mã & Thuộc tính & Tiêu chí có thể kiểm tra & Minh chứng cần có\\\midrule\endfirsthead
\toprule Mã & Thuộc tính & Tiêu chí có thể kiểm tra & Minh chứng cần có\\\midrule\endhead
NFR01 & Auth/ownership & Chưa auth: 401. Hai user không xem/sửa/xóa/link chéo ở mọi API dữ liệu riêng. & HTTP integration bằng A/B và DB assert.\\
NFR02 & Bảo mật web & Theo DEC07: cookie HttpOnly và Secure trên HTTPS; phiên lưu DB, kiểm tra hết hạn và thu hồi; CSRF bảo vệ mọi thao tác thay đổi dữ liệu. & AT16; kiểm thuộc tính cookie, CSRF, đăng xuất, hết hạn và đổi mật khẩu.\\
NFR03 & Toàn vẹn & FK/unique/check, rollback, idempotency và 409 version; không dữ liệu một phần. & PostgreSQL riêng; fault injection và concurrent writes.\\
NFR04 & Khả dụng UI & 360/768/1440px; keyboard/focus/label; trạng thái có chữ; không tràn toàn trang. & Checklist trình duyệt, ảnh và ghi nhận thao tác.\\
NFR05 & Hiệu năng & 10.000 hoạt động, 10 user đồng thời; p95 danh sách/báo cáo không quá 1 giây sau warmup. & Script tải, dữ liệu seed, raw samples, máy và DB/network.\\
NFR06 & Nhắc / phục hồi & Khi server bình thường, tạo nhắc thời gian trong 120 giây; tab active nhận trong 60 giây tiếp; restart không trùng/mất. & Controlled clock, worker logs và DB unique assertions.\\
NFR07 & Bảo trì & Các tầng tuân thủ hướng phụ thuộc; quy tắc nghiệp vụ có thể kiểm thử độc lập; quy trình build và kiểm tra được thống nhất. & Kết quả kiểm tra kiến trúc và các bước build, test.\\
NFR08 & Tái lập / bền vững & Checkout sạch, migration rõ ràng, restart giữ dữ liệu, restore trên DB riêng. & Runbook, log cài đặt/build và backup/restore drill.\\
NFR09 & Độ đúng & Cuối tháng/năm nhuận/VN date, tiền/ODO, ngày biên và sửa lịch sử cho kết quả đúng. & Test vectors có expected tính tay.\\
NFR10 & Quan sát / lỗi & ProblemDetails có code/traceId trong phản hồi lỗi; log không secret; ready phản ánh DB, worker hồi phục vòng sau. & Gây lỗi có kiểm soát, đối chiếu response/log/health.\\\bottomrule
\end{longtable}\end{smalltable}
